How can a computer, which rounds every number it touches, be trusted to solve a system of a million equations? How
fast can any method minimise a convex function, and why does a little momentum make it so much faster? What can a
sample of data tell us, with what confidence, and why does a machine that has seen finitely many examples predict the
next one?

This book proves the answers. It covers backward error analysis and the $QR$ algorithm, Krylov methods and random
sketches, convex duality, gradient descent with Nesterov's acceleration and its matching lower bound, compressed
sensing, interior point methods, online learning and bandits, estimation and the James--Stein paradox, tests, the false
discovery rate and the bootstrap, Bayesian inference and Markov chain Monte Carlo, the Vapnik--Chervonenkis theory of
generalisation, boosting, kernels and neural networks.

Every theorem is proved in full, every exercise has a hint and a reasoned solution, and every chapter comes with a
Python notebook that implements its methods and checks its numbers. Written for the pleasure of seeing why things work.
